\documentclass[12pt]{article}

\usepackage{amsmath,amssymb}
\usepackage{geometry}
\usepackage{graphicx}
\geometry{margin=1in}

\title{Lecture Scribe \\ CSE400: Fundamentals of Probability in Computing \\ Lecture 18}
\author{Topic: Marginal PMF, Correlation, Covariance, Joint Gaussian Random Variables}
\date{}

\begin{document}

\maketitle

\section{Marginal PMF}

\subsection{Definition}

Let $X$ and $Y$ be discrete random variables.

The \textbf{joint probability mass function} is defined as

\[
p_{XY}(x,y) = P(X=x, Y=y)
\]

The \textbf{marginal PMF} of $X$ is obtained by summing the joint PMF over all possible values of $Y$:

\[
p_X(x) = \sum_y p_{XY}(x,y)
\]

Similarly, the marginal PMF of $Y$ is

\[
p_Y(y) = \sum_x p_{XY}(x,y)
\]

\textbf{Key Idea:}

The marginal PMF of one variable is obtained by summing the joint PMF over all possible values of the other variable.

\subsection{Marginal PMF from a Joint PMF Table}

If the joint PMF is given in tabular form

\[
\begin{array}{c|ccc}
p_{XY}(x,y) & y_1 & y_2 & y_3 \\
\hline
x_1 & p_{11} & p_{12} & p_{13} \\
x_2 & p_{21} & p_{22} & p_{23} \\
x_3 & p_{31} & p_{32} & p_{33}
\end{array}
\]

Then

\[
p_X(x_i) = \sum_j p_{ij}
\]

\[
p_Y(y_j) = \sum_i p_{ij}
\]

and

\[
\sum_i p_X(x_i) = 1, \qquad \sum_j p_Y(y_j) = 1
\]

\subsection{Example}

Let the joint PMF be

\[
\begin{array}{c|cc}
p_{XY}(x,y) & y=0 & y=1 \\
\hline
x=0 & \frac18 & \frac14 \\
x=1 & \frac38 & \frac14
\end{array}
\]

\textbf{Computation}

\[
p_X(0) = \frac18 + \frac14 = \frac38
\]

\[
p_X(1) = \frac38 + \frac14 = \frac58
\]

\[
p_Y(0) = \frac18 + \frac38 = \frac12
\]

\[
p_Y(1) = \frac14 + \frac14 = \frac12
\]

\textbf{Conclusion}

\[
p_X(x) =
\begin{cases}
\frac38 & x=0 \\
\frac58 & x=1
\end{cases}
\]

\[
p_Y(y) =
\begin{cases}
\frac12 & y=0 \\
\frac12 & y=1
\end{cases}
\]

\section{Correlation}

\subsection{Definition}

A basic way to measure joint behavior between two variables is through the \textbf{product moment}

\[
R_{XY} = E[XY]
\]

It represents the average value of the product $XY$.

\textbf{Properties}

\begin{itemize}
\item $R_{XY}$ depends on the scale of $X$ and $Y$
\item It is computed about the origin
\item It mixes together mean values, spread, and dependence
\end{itemize}

If

\[
R_{XY} = 0
\]

then the variables are orthogonal about the origin.

To remove the effect of the means, variables are centered, leading to covariance.

\section{Covariance}

\subsection{Definition}

Covariance measures how two random variables vary together.

\[
\text{Cov}(X,Y) = E[(X-\mu_X)(Y-\mu_Y)]
\]

where

\[
\mu_X = E[X], \qquad \mu_Y = E[Y]
\]

\subsection{Algebraic Form}

\[
\text{Cov}(X,Y) = E[(X-\mu_X)(Y-\mu_Y)]
\]

Expanding,

\[
= E[XY - X\mu_Y - \mu_X Y + \mu_X \mu_Y]
\]

\[
= E[XY] - \mu_YE[X] - \mu_XE[Y] + \mu_X\mu_Y
\]

Since

\[
\mu_X = E[X], \quad \mu_Y = E[Y]
\]

we obtain

\[
\boxed{\text{Cov}(X,Y) = E[XY] - E[X]E[Y]}
\]

\subsection{Interpretation}

\begin{itemize}
\item $\text{Cov}(X,Y) > 0$ : positive covariance
\item $\text{Cov}(X,Y) < 0$ : negative covariance
\item $\text{Cov}(X,Y) = 0$ : zero covariance
\end{itemize}

Covariance indicates the direction of joint linear variation between $X$ and $Y$.

However, its magnitude depends on the units of the variables.

\subsection{Properties}

\textbf{1. Symmetry}

\[
\text{Cov}(X,Y) = \text{Cov}(Y,X)
\]

\textbf{2. Variance as a Special Case}

\[
\text{Cov}(X,X) = \text{Var}(X)
\]

\textbf{3. Linearity with Respect to Constants}

\[
\text{Cov}(aX+b,cY+d) = ac \, \text{Cov}(X,Y)
\]

\textbf{4. Independence}

If $X$ and $Y$ are independent,

\[
\text{Cov}(X,Y) = 0
\]

Zero covariance implies uncorrelated variables but does not necessarily imply independence.

\section{Example}

If

\[
Y = -2X + 3
\]

find $\text{Cov}(X,Y)$.

\textbf{Solution}

Using the covariance formula

\[
\text{Cov}(X,Y) = E[XY] - E[X]E[Y]
\]

Since

\[
Y = -2X + 3
\]

\[
XY = X(-2X+3) = -2X^2 + 3X
\]

\[
E[Y] = E[-2X+3] = -2E[X] + 3
\]

Substituting

\[
\text{Cov}(X,Y)
= E[-2X^2 + 3X] - E[X](-2E[X]+3)
\]

\[
= -2E[X^2] + 3E[X] + 2(E[X])^2 - 3E[X]
\]

\[
= -2E[X^2] + 2(E[X])^2
\]

Using

\[
\text{Var}(X) = E[X^2] - (E[X])^2
\]

we obtain

\[
\boxed{\text{Cov}(X,Y) = -2 \, \text{Var}(X)}
\]

\section{Pearson Correlation Coefficient}

The Pearson correlation coefficient measures the strength and direction of the linear relationship between $X$ and $Y$.

\[
\rho_{XY} =
\frac{\text{Cov}(X,Y)}{\sqrt{\text{Var}(X)\text{Var}(Y)}}
\]

or

\[
\rho_{XY} = \frac{\text{Cov}(X,Y)}{\sigma_X \sigma_Y}
\]

\textbf{Properties}

\begin{itemize}
\item It is the normalized version of covariance
\item It is unit-free
\item It measures both direction and strength
\end{itemize}

\[
-1 \le \rho_{XY} \le 1
\]

\[
\rho_{XY} > 0 \Rightarrow \text{positive linear trend}
\]

\[
\rho_{XY} < 0 \Rightarrow \text{negative linear trend}
\]

\[
\rho_{XY} = 0 \Rightarrow \text{no linear correlation}
\]

\section{Example 2}

Let the joint PDF be

\[
f_{XY}(x,y)=
\begin{cases}
\frac{xy}{96}, & 0<x<4, \; 1<y<5 \\
0 & \text{otherwise}
\end{cases}
\]

Compute

\begin{enumerate}
\item $E[X]$
\item $E[Y]$
\item $E[XY]$
\item $E[2X+3Y]$
\item $\text{Var}(X)$
\item $\text{Var}(Y)$
\item $\text{Cov}(X,Y)$
\item $\rho_{XY}$
\end{enumerate}

\subsection{Computation}

\[
E[X] = \int_0^4 x f_X(x) dx
= \frac18 \int_0^4 x^2 dx
= \frac83
\]

\[
E[Y] = \int_1^5 y f_Y(y) dy
= \frac1{12} \int_1^5 y^2 dy
= \frac{31}{9}
\]

\[
E[XY] = E[X]E[Y]
= \frac83 \cdot \frac{31}{9}
= \frac{248}{27}
\]

\[
E[2X+3Y] = 2E[X] + 3E[Y]
= \frac{47}{3}
\]

\[
\text{Var}(X) = \frac89
\]

\[
\text{Var}(Y) = \frac{92}{81}
\]

\[
\text{Cov}(X,Y) = 0
\]

\[
\rho_{XY} = 0
\]

\section{Jointly Gaussian Random Variables}

\subsection{Definition}

A pair $(X,Y)$ is jointly Gaussian if the joint density has the bivariate Gaussian form and the marginals are Gaussian.

\subsection{Joint Gaussian PDF}

\[
f_{XY}(x,y) =
\frac{1}{2\pi\sigma_X\sigma_Y\sqrt{1-\rho^2}}
\exp
\left(
-\frac{1}{2(1-\rho^2)}
\left[
\left(\frac{x-\mu_X}{\sigma_X}\right)^2
+
\left(\frac{y-\mu_Y}{\sigma_Y}\right)^2
-
2\rho
\left(\frac{x-\mu_X}{\sigma_X}\right)
\left(\frac{y-\mu_Y}{\sigma_Y}\right)
\right]
\right)
\]

where

\begin{itemize}
\item $\mu_X, \mu_Y$ are the means
\item $\sigma_X, \sigma_Y$ are the standard deviations
\item $\rho$ is the Pearson correlation coefficient
\end{itemize}

\end{document}